\documentclass[11pt]{article}
\usepackage[margin=1in]{geometry}
\usepackage{xcolor}
\usepackage{tcolorbox}
\usepackage{hyperref}
\usepackage{enumitem}
\usepackage{listings}

% Define colors
\definecolor{osaorange}{RGB}{220,115,38}
\definecolor{aiblue}{RGB}{0,83,155}

\title{\textbf{Group 2: Pima Diabetes Study}\\
\large{Dataset Guidance}}
\author{}
\date{}

\begin{document}

\maketitle

\section*{Dataset Overview}

\textbf{File:} \texttt{pima\_diabetes.csv}

\textbf{Sample Size:} n = 768 Pima Indian women

\textbf{Research Context:} Cross-sectional study examining risk factors for diabetes. Multiple health measures collected to identify which factors are associated with diabetes diagnosis.

\subsection*{Variables}

\begin{itemize}
\item \texttt{glu}: Plasma glucose concentration (mg/dL)
\item \texttt{bmi}: Body mass index (kg/m²)
\item \texttt{age}: Age (years)
\item \texttt{type}: Diabetes diagnosis (No or Yes)
\item Additional variables may include: blood pressure, pregnancies, skin fold thickness
\end{itemize}

\subsection*{Loading the Data in R}

\begin{lstlisting}[language=R, basicstyle=\small\ttfamily, breaklines=true]
# Load the dataset
pima <- read.csv("pima_diabetes.csv")

# View structure
str(pima)
head(pima)

# Summary by diabetes status
table(pima$type)
tapply(pima$glu, pima$type, summary)
tapply(pima$bmi, pima$type, summary)
\end{lstlisting}

\section*{Research Question to Consider}

\textbf{``Which health factors are associated with diabetes?''}

Alternatively: ``Do glucose levels, BMI, and age differ between women with and without diabetes?''

\section*{Key Decision Point}

\begin{tcolorbox}[colback=osaorange!10,colframe=osaorange,boxrule=1pt,arc=2pt]
\textbf{The Challenge:} You have a \textbf{binary outcome} (diabetes: Yes/No) and several continuous predictors (glucose, BMI, age). How do you analyze this with methods covered in this course?

\textbf{Questions to explore with AI:}
\begin{itemize}
\item What methods are commonly used for binary outcomes?
\item Can you compare groups (Yes vs. No) on each predictor separately?
\item What are the trade-offs between simple and complex approaches?
\end{itemize}

\textbf{Important note:} AI tools will likely suggest advanced methods for binary outcomes that you haven't learned yet. Consider whether simpler approaches can still answer your research question.
\end{tcolorbox}

\section*{Things to Think About}

\noindent\textbf{Exploratory Analysis:}
\begin{itemize}
\item Create visualizations comparing diabetes groups (boxplots, histograms)
\item Calculate summary statistics by diabetes status
\item Consider: Which variables show the biggest differences between groups?
\end{itemize}

\noindent\textbf{Statistical Methods:}
\begin{itemize}
\item Review methods for comparing groups
\item Consider: How do you compare two groups on a continuous variable?
\item Consider: What if you wanted to categorize continuous variables?
\end{itemize}

\noindent\textbf{Multiple Predictors:}
\begin{itemize}
\item You have several potential risk factors (glucose, BMI, age, etc.)
\item Should you analyze them one at a time or together?
\item What methods from this course allow you to examine each factor?
\end{itemize}

\section*{AI Consultation Strategy}

\begin{tcolorbox}[colback=aiblue!10,colframe=aiblue,boxrule=1pt,arc=2pt]
\textbf{Very Important:} When you ask AI about analyzing binary outcomes (Yes/No), it will almost certainly suggest methods beyond introductory biostatistics.

\textbf{Prompts to try:}
\begin{itemize}
\item ``How should I analyze which factors are associated with diabetes?''
\item ``I have a binary outcome and continuous predictors - what methods can I use?''
\item ``I haven't learned [method AI suggested]. Can I answer my question with methods covered in this course?''
\item ``What's the difference between analyzing predictors one at a time vs. together?''
\end{itemize}

\textbf{Watch for:}
\begin{itemize}
\item Tools suggesting methods you haven't encountered (especially for binary outcomes)
\item Different recommendations from different AI tools
\item Trade-offs between simple (univariate) and complex (multivariate) approaches
\end{itemize}

\textbf{Critical thinking:} Just because AI suggests an advanced method doesn't mean simpler methods can't answer your question. Can you identify risk factors using methods you've learned?
\end{tcolorbox}

\section*{Questions for Your Analysis}

As you work through this dataset, consider:

\begin{enumerate}
\item What research question are you answering?
\item Which statistical method(s) from this course can address this question?
\item Are you analyzing predictors individually or together? Why?
\item What are the limitations of your chosen approach?
\item How will you justify your methodological choice when AI suggested something different?
\end{enumerate}

\section*{Interpretation Considerations}

If you find significant differences between diabetes groups:
\begin{itemize}
\item What does this tell you about risk factors?
\item How would you interpret findings for clinical practice?
\item What can you conclude from examining factors one at a time?
\item What questions remain that would require more advanced methods?
\end{itemize}

\section*{Resources}

\begin{itemize}
\item Week 9 drop-in session (Wed 12-1:20pm, 153 Milam)
\item Textbook: Chapter 11 (Two-Sample Tests), Chapter 15 (Chi-Square)
\item Your team!
\item Canvas discussion board
\end{itemize}

\vspace{0.5cm}

\begin{tcolorbox}[colback=gray!10,colframe=gray!50,boxrule=1pt,arc=2pt]
\textbf{Remember:} Simple methods that you understand are often better than complex methods you don't. The goal is to:
\begin{itemize}
\item Answer your research question with appropriate methods from this course
\item Understand and justify your analytical choices
\item Critically evaluate what AI tools suggested and why
\item Recognize when advanced methods exist but aren't necessary
\end{itemize}
\end{tcolorbox}

\end{document}
